\documentclass[UTF8,12pt]{ctexart}

\usepackage[a4paper,margin=2.3cm]{geometry}
\usepackage{amsmath,amssymb,mathtools,bm}
\usepackage{booktabs,longtable,array}
\usepackage{enumitem}
\usepackage{hyperref}
\usepackage{xcolor}

\hypersetup{
  colorlinks=true,
  linkcolor=blue,
  citecolor=blue,
  urlcolor=blue
}

\setlist[itemize]{leftmargin=2em}
\setlist[enumerate]{leftmargin=2em}

\newcommand{\R}{\mathbb{R}}
\newcommand{\E}{\mathbb{E}}
\newcommand{\1}{\mathbf{1}}
\newcommand{\softmax}{\operatorname{softmax}}
\newcommand{\argmax}{\operatorname*{arg\,max}}
\newcommand{\TopK}{\operatorname{TopK}}
\newcommand{\LN}{\operatorname{LN}}
\newcommand{\RMSNorm}{\operatorname{RMSNorm}}
\newcommand{\Attn}{\operatorname{Attention}}
\newcommand{\MHA}{\operatorname{MHA}}
\newcommand{\FFN}{\operatorname{FFN}}
\newcommand{\GELU}{\operatorname{GELU}}
\newcommand{\SiLU}{\operatorname{SiLU}}
\newcommand{\SwiGLU}{\operatorname{SwiGLU}}
\newcommand{\KL}{D_{\mathrm{KL}}}
\newcommand{\NLL}{\operatorname{NLL}}
\newcommand{\PPL}{\operatorname{PPL}}
\newcommand{\diag}{\operatorname{diag}}
\newcommand{\clip}{\operatorname{clip}}
\newcommand{\sgn}{\operatorname{sgn}}

\title{LLM 数学公式参考：从 Transformer 到现代大模型推理系统}
\author{miniLLM 学习笔记}
\date{\today}

\begin{document}
\maketitle
\tableofcontents
\newpage

\section{怎么读这份笔记}

这份笔记按 LLM 的真实数据流组织公式：

\[
\text{文本}
\rightarrow \text{Tokenizer}
\rightarrow \text{Token IDs}
\rightarrow \text{Embedding}
\rightarrow \text{Transformer Blocks}
\rightarrow \text{Logits}
\rightarrow
\begin{cases}
\text{训练：Loss + Optimizer}\\
\text{推理：Sampling + KV Cache}
\end{cases}
\]

其中最核心的数学链条是：

\[
p_\theta(x_{1:T})
=\prod_{t=1}^{T}p_\theta(x_t\mid x_{<t})
\quad\Longrightarrow\quad
\mathcal{L}_{\mathrm{LM}}
=-\sum_{t=1}^{T}\log p_\theta(x_t\mid x_{<t})
\quad\Longrightarrow\quad
\nabla_\theta \mathcal{L}
\quad\Longrightarrow\quad
\theta \leftarrow \theta-\eta \cdot \widehat{\nabla_\theta \mathcal{L}}.
\]

Transformer 只是实现条件概率 $p_\theta(x_t\mid x_{<t})$ 的函数族；vLLM/SGLang 这类推理系统则负责高效执行这个函数族。

\section{符号约定}

\begin{longtable}{lll}
\toprule
符号 & 形状/含义 & 说明\\
\midrule
$V$ & 词表大小 & vocabulary size\\
$T$ & 序列长度 & sequence length\\
$B$ & batch size & 批大小\\
$d_{\mathrm{model}}$ & hidden size & 模型宽度\\
$H$ & attention heads & query heads 数\\
$H_{\mathrm{kv}}$ & KV heads & MHA: $H_{\mathrm{kv}}=H$；MQA: $1$；GQA: $G$\\
$d_h$ & head dim & 通常 $d_{\mathrm{model}}=H d_h$\\
$L$ & Transformer 层数 & number of layers\\
$x_t$ & 第 $t$ 个 token id & $x_t\in\{1,\dots,V\}$\\
$\bm{h}_t^\ell$ & 第 $\ell$ 层第 $t$ 个 hidden state & $\bm{h}_t^\ell\in\R^{d_{\mathrm{model}}}$\\
$\theta$ & 模型参数 & all trainable parameters\\
\bottomrule
\end{longtable}

\section{语言建模的概率基础}

\subsection{自回归分解}

LLM 的预训练核心是估计 token 序列的联合概率。根据链式法则：

\begin{equation}
p_\theta(x_{1:T})
=\prod_{t=1}^{T}p_\theta(x_t\mid x_{1:t-1}).
\label{eq:ar-factorization}
\end{equation}

GPT 类模型使用这个自回归分解训练和生成。来源：GPT/GPT-2 系列语言模型目标 \cite{radford2019gpt2,brown2020gpt3}，Transformer decoder 结构来源于 \cite{vaswani2017attention}。

\subsection{最大似然、交叉熵和困惑度}

对数据集 $\mathcal{D}$ 最大化似然：

\begin{equation}
\theta^\star
=\argmax_\theta
\sum_{x_{1:T}\in\mathcal{D}}\log p_\theta(x_{1:T}).
\end{equation}

等价于最小化 negative log-likelihood：

\begin{equation}
\mathcal{L}_{\mathrm{NLL}}(\theta)
=-\sum_{x_{1:T}\in\mathcal{D}}\sum_{t=1}^{T}
\log p_\theta(x_t\mid x_{<t}).
\label{eq:nll}
\end{equation}

若对 token 平均：

\begin{equation}
\bar{\mathcal{L}}_{\mathrm{NLL}}
=-\frac{1}{N_{\mathrm{tok}}}
\sum_{n=1}^{N_{\mathrm{tok}}}
\log p_\theta(x_n\mid x_{<n}).
\end{equation}

困惑度：

\begin{equation}
\PPL=\exp\left(\bar{\mathcal{L}}_{\mathrm{NLL}}\right).
\end{equation}

解释：

\begin{itemize}
  \item NLL 越低，模型给真实下一个 token 的概率越高。
  \item PPL 可理解为模型平均在多少个等概率候选中困惑；越低越好。
\end{itemize}

\subsection{Masked LM 目标}

BERT 使用 masked language modeling，不是纯自回归：

\begin{equation}
\mathcal{L}_{\mathrm{MLM}}
=-\sum_{i\in\mathcal{M}}
\log p_\theta(x_i\mid x_{\setminus \mathcal{M}}),
\label{eq:mlm}
\end{equation}

其中 $\mathcal{M}$ 是被 mask 的 token 位置集合。来源：BERT \cite{devlin2018bert}。

关系：

\begin{itemize}
  \item GPT/LLM 推理常用公式是式 \eqref{eq:ar-factorization}。
  \item BERT 更适合理解/编码任务；GPT 更适合逐 token 生成。
\end{itemize}

\section{Tokenizer 与离散 token}

\subsection{BPE 合并目标}

BPE 每一步选择语料中出现最频繁的相邻 symbol pair 合并：

\begin{equation}
(a^\star,b^\star)
=\argmax_{(a,b)}
\operatorname{count}_{\mathcal{D}}(a,b).
\end{equation}

来源：subword BPE for NMT \cite{sennrich2015bpe}。

\subsection{Unigram tokenizer}

SentencePiece 的 unigram 模型把一句话 $\bm{x}$ 的分词看成潜变量 $\bm{z}$：

\begin{equation}
p(\bm{x})
=\sum_{\bm{z}\in \mathcal{Z}(\bm{x})}
\prod_{u\in \bm{z}}p(u).
\end{equation}

训练时最大化语料 likelihood：

\begin{equation}
\mathcal{L}_{\mathrm{tok}}
=\sum_{\bm{x}\in\mathcal{D}}\log p(\bm{x}).
\end{equation}

来源：SentencePiece / subword regularization \cite{kudo2018sentencepiece,kudo2018subword}。

关系：

\[
\text{Tokenizer}
\quad\Longrightarrow\quad
x_t\in\{1,\dots,V\}
\quad\Longrightarrow\quad
\text{Embedding lookup}.
\]

\section{Embedding、输出投影与 tied weights}

\subsection{Token embedding}

设 token embedding 矩阵：

\begin{equation}
E\in\R^{V\times d_{\mathrm{model}}}.
\end{equation}

token id $x_t$ 对应 one-hot 向量 $\bm{e}_{x_t}\in\R^V$，embedding 为：

\begin{equation}
\bm{u}_t=E^\top \bm{e}_{x_t}
=E[x_t]\in\R^{d_{\mathrm{model}}}.
\end{equation}

\subsection{输出 logits}

最后一层 hidden state $\bm{h}_t^L$ 投影到词表：

\begin{equation}
\bm{z}_t=W_U \bm{h}_t^L+\bm{b}_U,
\quad
W_U\in\R^{V\times d_{\mathrm{model}}}.
\end{equation}

概率分布：

\begin{equation}
p_\theta(x_{t+1}=i\mid x_{\le t})
=\softmax(\bm{z}_t)_i
=\frac{\exp(z_{t,i})}{\sum_{j=1}^{V}\exp(z_{t,j})}.
\label{eq:softmax}
\end{equation}

\subsection{Weight tying}

很多 LLM 使用输入 embedding 和输出投影共享权重：

\begin{equation}
W_U = E.
\end{equation}

这样 logits 可写成：

\begin{equation}
z_{t,i}=E[i]^\top \bm{h}_t^L+b_i.
\end{equation}

解释：输出 token $i$ 的 logit 是 hidden state 与 token embedding 的相似度。

\section{Transformer Block 总览}

\subsection{Pre-LN Transformer block}

现代 decoder-only LLM 多采用 Pre-LN 或 RMSNorm 变体：

\begin{align}
\tilde{\bm{h}}^\ell
&=\bm{h}^{\ell}
+\MHA(\operatorname{Norm}(\bm{h}^{\ell})),\\
\bm{h}^{\ell+1}
&=\tilde{\bm{h}}^\ell
+\FFN(\operatorname{Norm}(\tilde{\bm{h}}^\ell)).
\label{eq:preln-block}
\end{align}

Transformer 原始结构来源：\cite{vaswani2017attention}。Pre-LN 是后续稳定训练常见工程选择。

依赖关系：

\[
\text{Norm}
\rightarrow
\text{Attention}
\rightarrow
\text{Residual}
\rightarrow
\text{Norm}
\rightarrow
\text{FFN}
\rightarrow
\text{Residual}.
\]

\section{LayerNorm 与 RMSNorm}

\subsection{LayerNorm}

对向量 $\bm{x}\in\R^d$：

\begin{align}
\mu &= \frac{1}{d}\sum_{i=1}^{d}x_i,\\
\sigma^2 &= \frac{1}{d}\sum_{i=1}^{d}(x_i-\mu)^2,\\
\LN(\bm{x})
&=\bm{\gamma}\odot
\frac{\bm{x}-\mu}{\sqrt{\sigma^2+\epsilon}}
+\bm{\beta}.
\label{eq:layernorm}
\end{align}

来源：Layer Normalization \cite{ba2016layernorm}。

\subsection{RMSNorm}

RMSNorm 去掉均值中心化，只除以 root mean square：

\begin{align}
\operatorname{RMS}(\bm{x})
&=\sqrt{\frac{1}{d}\sum_{i=1}^{d}x_i^2+\epsilon},\\
\RMSNorm(\bm{x})
&=\bm{g}\odot \frac{\bm{x}}{\operatorname{RMS}(\bm{x})}.
\label{eq:rmsnorm}
\end{align}

来源：RMSNorm \cite{zhang2019rmsnorm}。

关系：

\begin{itemize}
  \item LayerNorm 控制均值和方差。
  \item RMSNorm 只控制尺度，计算更少，很多 LLaMA-like 模型使用 RMSNorm。
\end{itemize}

\section{Attention：从 QK 到 Softmax}

\subsection{Q、K、V 的线性投影}

对输入 hidden states $X\in\R^{T\times d_{\mathrm{model}}}$：

\begin{align}
Q &= XW_Q,\quad W_Q\in\R^{d_{\mathrm{model}}\times d_k},\\
K &= XW_K,\quad W_K\in\R^{d_{\mathrm{model}}\times d_k},\\
V &= XW_V,\quad W_V\in\R^{d_{\mathrm{model}}\times d_v}.
\end{align}

含义：

\begin{itemize}
  \item Query：当前位置想找什么信息。
  \item Key：每个上下文位置提供什么索引。
  \item Value：真正被聚合的信息。
\end{itemize}

\subsection{为什么是 $QK^\top$}

第 $i$ 个 query 与第 $j$ 个 key 的相似度：

\begin{equation}
s_{ij}=\bm{q}_i^\top \bm{k}_j.
\end{equation}

把所有位置一次性写成矩阵：

\begin{equation}
S=QK^\top,\quad S\in\R^{T\times T}.
\end{equation}

缩放点积注意力：

\begin{equation}
\Attn(Q,K,V)
=\softmax\left(\frac{QK^\top}{\sqrt{d_k}}\right)V.
\label{eq:scaled-dot-product-attention}
\end{equation}

来源：Transformer \cite{vaswani2017attention}。

为什么除以 $\sqrt{d_k}$：

若 $q_i,k_i$ 独立且均值 0、方差 1，则

\begin{equation}
\operatorname{Var}(\bm{q}^\top\bm{k})
=\operatorname{Var}\left(\sum_{r=1}^{d_k}q_r k_r\right)
=d_k.
\end{equation}

除以 $\sqrt{d_k}$ 后方差约为 1，避免 softmax 输入随维度变大而过大，导致梯度饱和。

\subsection{Causal mask}

Decoder-only LLM 不能看未来 token：

\begin{equation}
M_{ij}=
\begin{cases}
0, & j\le i,\\
-\infty, & j>i.
\end{cases}
\end{equation}

于是：

\begin{equation}
\Attn_{\mathrm{causal}}(Q,K,V)
=\softmax\left(\frac{QK^\top}{\sqrt{d_k}}+M\right)V.
\label{eq:causal-attention}
\end{equation}

\subsection{Softmax 的数值稳定形式}

直接计算 $\exp(z_i)$ 可能溢出，通常减去最大值：

\begin{equation}
\softmax(\bm{z})_i
=\frac{\exp(z_i-m)}{\sum_j \exp(z_j-m)},
\quad
m=\max_j z_j.
\label{eq:stable-softmax}
\end{equation}

这不改变结果，因为分子分母同除 $\exp(m)$。

\subsection{为什么 softmax/attention 内存受限}

完整 attention score 的元素数：

\begin{equation}
\#S = B\cdot H\cdot T^2.
\end{equation}

若用 fp16/bf16，每个元素 2 bytes，仅 score 矩阵显存读写量约为：

\begin{equation}
M_S \approx 2\cdot BHT^2\ \mathrm{bytes}.
\end{equation}

实际还会有 softmax 概率矩阵、dropout mask、反向传播缓存等。因此长序列训练中，attention 往往受 HBM 带宽和中间矩阵 materialization 限制。

FlashAttention 的思想是不要显式 materialize 完整 $T\times T$ attention matrix，而是分块做 online softmax。来源：FlashAttention \cite{dao2022flashattention}、FlashAttention-2 \cite{dao2023flashattention2}。

\subsection{FlashAttention 的 online softmax 递推}

对一行 attention score 分块处理。维护：

\begin{align}
m &= \max_j s_j,\\
\ell &= \sum_j \exp(s_j-m),\\
\bm{n} &= \sum_j \exp(s_j-m)\bm{v}_j,\\
\bm{o} &= \frac{\bm{n}}{\ell}.
\end{align}

当新块 $B$ 到来：

\begin{align}
m' &= \max(m,\max_{j\in B}s_j),\\
\ell' &= e^{m-m'}\ell+\sum_{j\in B}e^{s_j-m'},\\
\bm{n}' &= e^{m-m'}\bm{n}+\sum_{j\in B}e^{s_j-m'}\bm{v}_j,\\
\bm{o}' &= \frac{\bm{n}'}{\ell'}.
\end{align}

这个递推保证结果等价于完整 softmax，但显著减少 HBM 读写。来源：\cite{dao2022flashattention}。

\section{Multi-Head、MQA 与 GQA}

\subsection{Multi-Head Attention}

第 $h$ 个 head：

\begin{equation}
\operatorname{head}_h
=\Attn(XW_Q^h,XW_K^h,XW_V^h).
\end{equation}

拼接所有 heads：

\begin{equation}
\MHA(X)
=\operatorname{Concat}(\operatorname{head}_1,\dots,\operatorname{head}_H)W_O.
\label{eq:mha}
\end{equation}

来源：\cite{vaswani2017attention}。

\subsection{Multi-Query Attention}

MQA 让多个 query heads 共享同一组 K/V：

\begin{align}
Q_h &= XW_Q^h,\quad h=1,\dots,H,\\
K &= XW_K,\quad V=XW_V,\\
\operatorname{head}_h &= \Attn(Q_h,K,V).
\end{align}

KV cache 从 $H$ 组变成 1 组，显著降低解码时显存带宽和 KV cache 大小。来源：Multi-Query Attention \cite{shazeer2019mqa}。

\subsection{Grouped-Query Attention}

GQA 是 MHA 和 MQA 的折中。设 $G=H_{\mathrm{kv}}$ 个 KV groups，query head $h$ 使用 group：

\begin{equation}
g(h)=\left\lfloor \frac{hG}{H}\right\rfloor.
\end{equation}

则：

\begin{equation}
\operatorname{head}_h
=\Attn(Q_h,K_{g(h)},V_{g(h)}).
\end{equation}

来源：GQA \cite{ainslie2023gqa}。

关系：

\[
\text{MHA}:H_{\mathrm{kv}}=H,\quad
\text{GQA}:1<H_{\mathrm{kv}}<H,\quad
\text{MQA}:H_{\mathrm{kv}}=1.
\]

\section{位置编码}

\subsection{Sinusoidal positional encoding}

Transformer 原始位置编码：

\begin{align}
PE(pos,2i)
&=\sin\left(\frac{pos}{10000^{2i/d_{\mathrm{model}}}}\right),\\
PE(pos,2i+1)
&=\cos\left(\frac{pos}{10000^{2i/d_{\mathrm{model}}}}\right).
\end{align}

输入变为：

\begin{equation}
\bm{h}_t^0=E[x_t]+PE(t).
\end{equation}

来源：\cite{vaswani2017attention}。

\subsection{RoPE}

RoPE 对 query/key 做位置相关旋转。对二维子空间：

\begin{equation}
R_{\theta,m}
=
\begin{bmatrix}
\cos(m\theta)&-\sin(m\theta)\\
\sin(m\theta)&\cos(m\theta)
\end{bmatrix}.
\end{equation}

位置 $m$ 的 query 和位置 $n$ 的 key：

\begin{equation}
\tilde{\bm{q}}_m=R_m\bm{q}_m,\quad
\tilde{\bm{k}}_n=R_n\bm{k}_n.
\end{equation}

注意力分数：

\begin{equation}
\tilde{\bm{q}}_m^\top \tilde{\bm{k}}_n
=\bm{q}_m^\top R_m^\top R_n\bm{k}_n
=\bm{q}_m^\top R_{n-m}\bm{k}_n.
\label{eq:rope-relative}
\end{equation}

所以 RoPE 天然把相对位置 $n-m$ 注入 attention score。来源：RoFormer/RoPE \cite{su2021rope}。

\subsection{ALiBi}

ALiBi 直接给 attention score 加线性距离偏置：

\begin{equation}
s_{h,i,j}
=\frac{\bm{q}_{h,i}^\top \bm{k}_{h,j}}{\sqrt{d_h}}
-m_h(i-j),
\quad j\le i,
\end{equation}

其中 $m_h$ 是 head-specific slope。来源：ALiBi \cite{press2021alibi}。

关系：

\begin{itemize}
  \item Sinusoidal：位置向量加到输入 hidden。
  \item RoPE：位置旋转进 Q/K，使 attention score 依赖相对距离。
  \item ALiBi：直接对 attention score 加距离惩罚。
\end{itemize}

\section{FFN、GELU 与 SwiGLU}

\subsection{标准 FFN}

Transformer 每层还有 position-wise FFN：

\begin{equation}
\FFN(\bm{x})
=W_2\,\phi(W_1\bm{x}+\bm{b}_1)+\bm{b}_2.
\end{equation}

原始 Transformer 使用 ReLU。来源：\cite{vaswani2017attention}。

\subsection{GELU}

GELU：

\begin{equation}
\GELU(x)
=x\Phi(x),
\end{equation}

其中 $\Phi$ 是标准正态分布 CDF。常用近似：

\begin{equation}
\GELU(x)
\approx
0.5x\left(1+\tanh\left[\sqrt{\frac{2}{\pi}}\left(x+0.044715x^3\right)\right]\right).
\end{equation}

\subsection{SwiGLU}

现代 LLaMA-like 模型常用 SwiGLU：

\begin{align}
\SiLU(x)&=x\sigma(x),\\
\SwiGLU(\bm{x})
&=\left(\SiLU(\bm{x}W_1)\odot \bm{x}W_3\right)W_2.
\label{eq:swiglu}
\end{align}

来源：GLU variants \cite{shazeer2020glu}。

解释：$W_3$ 分支像 gate，控制 $W_1$ 分支的信息通过量。

\section{Softmax 交叉熵的梯度}

设 logits $\bm{z}\in\R^V$，概率 $\bm{p}=\softmax(\bm{z})$，真实标签 one-hot 为 $\bm{y}$：

\begin{equation}
\mathcal{L}_{\mathrm{CE}}
=-\sum_{i=1}^{V}y_i\log p_i.
\end{equation}

先求：

\begin{equation}
\log p_k
=z_k-\log\sum_j e^{z_j}.
\end{equation}

对 $z_i$ 求导：

\begin{equation}
\frac{\partial \log p_k}{\partial z_i}
=\1[i=k]-p_i.
\end{equation}

因此：

\begin{align}
\frac{\partial \mathcal{L}_{\mathrm{CE}}}{\partial z_i}
&=-\sum_k y_k\frac{\partial \log p_k}{\partial z_i}\\
&=-\sum_k y_k(\1[i=k]-p_i)\\
&=p_i-y_i.
\label{eq:ce-gradient}
\end{align}

这是 LLM 训练中最重要的梯度公式之一：

\[
\boxed{\nabla_{\bm{z}}\mathcal{L}_{\mathrm{CE}}=\bm{p}-\bm{y}}.
\]

关系：

\[
\text{logits}
\rightarrow
\softmax
\rightarrow
\text{cross entropy}
\rightarrow
\bm{p}-\bm{y}
\rightarrow
\text{backprop into Transformer}.
\]

\subsection{Label smoothing}

真实 one-hot $\bm{y}$ 可替换为：

\begin{equation}
\bm{y}'=(1-\epsilon)\bm{y}+\frac{\epsilon}{V}\bm{1}.
\end{equation}

则梯度变为：

\begin{equation}
\nabla_{\bm{z}}\mathcal{L}=\bm{p}-\bm{y}'.
\end{equation}

\section{优化器和训练动态}

\subsection{SGD}

最基础更新：

\begin{equation}
\theta_{t+1}=\theta_t-\eta_t\nabla_\theta \mathcal{L}(\theta_t).
\end{equation}

\subsection{Adam}

Adam 维护一阶、二阶矩估计：

\begin{align}
\bm{g}_t&=\nabla_\theta \mathcal{L}(\theta_t),\\
\bm{m}_t&=\beta_1\bm{m}_{t-1}+(1-\beta_1)\bm{g}_t,\\
\bm{v}_t&=\beta_2\bm{v}_{t-1}+(1-\beta_2)\bm{g}_t^2.
\end{align}

偏差修正：

\begin{equation}
\hat{\bm{m}}_t=\frac{\bm{m}_t}{1-\beta_1^t},
\quad
\hat{\bm{v}}_t=\frac{\bm{v}_t}{1-\beta_2^t}.
\end{equation}

更新：

\begin{equation}
\theta_{t+1}
=\theta_t-\eta_t\frac{\hat{\bm{m}}_t}{\sqrt{\hat{\bm{v}}_t}+\epsilon}.
\end{equation}

\subsection{AdamW}

AdamW 把 weight decay 从梯度里解耦：

\begin{equation}
\theta_{t+1}
=\theta_t
-\eta_t\frac{\hat{\bm{m}}_t}{\sqrt{\hat{\bm{v}}_t}+\epsilon}
-\eta_t\lambda \theta_t.
\label{eq:adamw}
\end{equation}

来源：AdamW \cite{loshchilov2017adamw}。

\subsection{Gradient clipping}

若全局梯度范数过大：

\begin{equation}
\bm{g}\leftarrow
\bm{g}\cdot
\min\left(1,\frac{c}{\|\bm{g}\|_2}\right).
\end{equation}

作用：防止梯度爆炸。

\subsection{Warmup + cosine decay}

常见学习率：

\begin{equation}
\eta(t)=
\begin{cases}
\eta_{\max}\frac{t}{T_{\mathrm{warmup}}}, & t<T_{\mathrm{warmup}},\\
\eta_{\min}
+\frac{1}{2}(\eta_{\max}-\eta_{\min})
\left(1+\cos\frac{\pi(t-T_{\mathrm{warmup}})}{T_{\mathrm{total}}-T_{\mathrm{warmup}}}\right),
& t\ge T_{\mathrm{warmup}}.
\end{cases}
\end{equation}

\section{Scaling Laws 与 Chinchilla}

\subsection{经验损失分解}

Scaling laws 常把 loss 写成模型大小 $N$、数据 token 数 $D$ 的函数：

\begin{equation}
L(N,D)
\approx L_\infty+\frac{A}{N^\alpha}+\frac{B}{D^\beta}.
\label{eq:scaling-laws}
\end{equation}

来源：Kaplan et al. scaling laws \cite{kaplan2020scaling}，Chinchilla \cite{hoffmann2022chinchilla}。

\subsection{训练计算量近似}

Transformer dense LM 的训练 FLOPs 常用近似：

\begin{equation}
C\approx 6ND,
\label{eq:training-compute}
\end{equation}

其中 $N$ 是非 embedding 参数量，$D$ 是训练 token 数。来源：Chinchilla \cite{hoffmann2022chinchilla}。

关系：

\[
C\ \text{固定}
\quad\Longrightarrow\quad
\text{选择 }N,D\text{ 的配比}
\quad\Longrightarrow\quad
\text{compute-optimal training}.
\]

直观结论：

\begin{itemize}
  \item 模型太大、数据太少：under-trained。
  \item 数据太多、模型太小：capacity 不足。
  \item Chinchilla 强调在固定 compute 下需要更多 tokens、更小但训练充分的模型。
\end{itemize}

\section{SFT、Reward Model、RLHF 和 DPO}

\subsection{Supervised Fine-Tuning}

给定 prompt $x$ 和参考回答 $y_{1:T}$：

\begin{equation}
\mathcal{L}_{\mathrm{SFT}}
=-\sum_{t=1}^{T}
\log \pi_\theta(y_t\mid x,y_{<t}).
\label{eq:sft}
\end{equation}

它和预训练 NLL 形式一样，只是数据变成 instruction-response。

\subsection{Reward model}

RLHF 中，人类偏好数据常为 $(x,y_w,y_l)$：$y_w$ 优于 $y_l$。Bradley--Terry 形式：

\begin{equation}
P(y_w\succ y_l\mid x)
=\sigma\left(r_\phi(x,y_w)-r_\phi(x,y_l)\right).
\end{equation}

reward model loss：

\begin{equation}
\mathcal{L}_{\mathrm{RM}}
=-\E_{(x,y_w,y_l)}
\log \sigma\left(r_\phi(x,y_w)-r_\phi(x,y_l)\right).
\label{eq:reward-model}
\end{equation}

来源：InstructGPT/RLHF \cite{ouyang2022instructgpt}。

\subsection{PPO for RLHF}

PPO 的 clipped objective：

\begin{align}
r_t(\theta)
&=\frac{\pi_\theta(a_t\mid s_t)}{\pi_{\theta_{\mathrm{old}}}(a_t\mid s_t)},\\
\mathcal{L}_{\mathrm{PPO}}
&=\E_t\left[
\min\left(
r_t(\theta)\hat{A}_t,
\clip(r_t(\theta),1-\epsilon,1+\epsilon)\hat{A}_t
\right)
\right].
\label{eq:ppo}
\end{align}

来源：PPO \cite{schulman2017ppo}。

LLM RLHF 通常还加 KL penalty，避免策略偏离 reference model：

\begin{equation}
\max_\theta
\E_{y\sim \pi_\theta(\cdot\mid x)}
\left[
r_\phi(x,y)
-\beta \KL\left(\pi_\theta(\cdot\mid x)\|\pi_{\mathrm{ref}}(\cdot\mid x)\right)
\right].
\label{eq:rlhf-kl}
\end{equation}

来源：InstructGPT \cite{ouyang2022instructgpt}。

\subsection{DPO}

DPO 直接从偏好对优化策略，不显式训练 reward model 做 RL rollouts：

\begin{align}
\mathcal{L}_{\mathrm{DPO}}(\theta)
=-\E_{(x,y_w,y_l)}
\log\sigma\Big(
\beta[
&\log\pi_\theta(y_w\mid x)-\log\pi_{\mathrm{ref}}(y_w\mid x)\\
-&\log\pi_\theta(y_l\mid x)+\log\pi_{\mathrm{ref}}(y_l\mid x)]
\Big).
\label{eq:dpo}
\end{align}

来源：DPO \cite{rafailov2023dpo}。

关系：

\[
\text{Pretrain}
\rightarrow
\text{SFT}
\rightarrow
\begin{cases}
\text{Reward Model}+\text{PPO}\\
\text{DPO}
\end{cases}
\rightarrow
\text{Aligned LLM}.
\]

\section{LoRA 与参数高效微调}

\subsection{低秩适配}

冻结原始权重 $W_0\in\R^{d_{\mathrm{out}}\times d_{\mathrm{in}}}$，只学习低秩增量：

\begin{equation}
W'=W_0+\Delta W,
\quad
\Delta W=\frac{\alpha}{r}BA,
\end{equation}

其中：

\begin{equation}
B\in\R^{d_{\mathrm{out}}\times r},
\quad
A\in\R^{r\times d_{\mathrm{in}}},
\quad
r\ll \min(d_{\mathrm{in}},d_{\mathrm{out}}).
\end{equation}

前向：

\begin{equation}
\bm{y}=W_0\bm{x}+\frac{\alpha}{r}BA\bm{x}.
\label{eq:lora}
\end{equation}

来源：LoRA \cite{hu2021lora}。

关系：

\[
\text{大矩阵全量更新}
\quad\Longrightarrow\quad
\text{低秩增量更新}
\quad\Longrightarrow\quad
\text{显著减少可训练参数和优化器状态}.
\]

\section{MoE：Dense 与 Sparse}

\subsection{Dense FFN}

Dense Transformer 每个 token 都经过同一个 FFN：

\begin{equation}
\bm{y}_t=\FFN(\bm{x}_t).
\end{equation}

每个 token 激活同一组参数。

\subsection{Sparse MoE}

有 $E$ 个 experts，router 输出：

\begin{equation}
\bm{p}_t=\softmax(W_r\bm{x}_t).
\end{equation}

选择 top-$k$ experts：

\begin{equation}
\mathcal{E}_t=\TopK(\bm{p}_t,k).
\end{equation}

输出：

\begin{equation}
\bm{y}_t
=\sum_{e\in\mathcal{E}_t}
p_{t,e}\operatorname{Expert}_e(\bm{x}_t).
\label{eq:moe}
\end{equation}

来源：GShard \cite{lepikhin2020gshard}、Switch Transformer \cite{fedus2021switch}。

\subsection{Switch Transformer top-1}

Switch Transformer 使用 top-1 routing：

\begin{equation}
e^\star_t=\argmax_e p_{t,e},
\quad
\bm{y}_t=p_{t,e^\star_t}\operatorname{Expert}_{e^\star_t}(\bm{x}_t).
\end{equation}

\subsection{负载均衡损失}

设 $f_e$ 是分配到 expert $e$ 的 token 比例，$P_e$ 是 router 给 expert $e$ 的平均概率，Switch 的辅助损失：

\begin{equation}
\mathcal{L}_{\mathrm{aux}}
=\alpha E\sum_{e=1}^{E} f_e P_e.
\label{eq:moe-aux}
\end{equation}

来源：Switch Transformer \cite{fedus2021switch}。

关系：

\begin{itemize}
  \item Dense：所有 token 都走相同参数，计算稳定但参数利用密集。
  \item Sparse MoE：每个 token 只激活少数 experts，总参数可很大，单 token 计算可控。
  \item Expert Parallelism：把 experts 分布到不同 GPU，需要 all-to-all token dispatch/combine。
\end{itemize}

\section{推理：从 logits 到 token}

\subsection{Temperature}

训练时使用 $\softmax(\bm{z})$，推理时可调 temperature $\tau$：

\begin{equation}
p_\tau(i)
=\frac{\exp(z_i/\tau)}{\sum_j \exp(z_j/\tau)}.
\label{eq:temperature}
\end{equation}

\begin{itemize}
  \item $\tau<1$：分布更尖锐，更保守。
  \item $\tau>1$：分布更平，更随机。
\end{itemize}

\subsection{Greedy decoding}

\begin{equation}
x_{t+1}=\argmax_i p_\theta(i\mid x_{\le t}).
\end{equation}

\subsection{Top-k sampling}

只保留概率最高的 $k$ 个 token：

\begin{equation}
S_k=\TopK(\bm{p},k),
\quad
\tilde{p}_i=
\begin{cases}
\frac{p_i}{\sum_{j\in S_k}p_j}, & i\in S_k,\\
0, & i\notin S_k.
\end{cases}
\end{equation}

\subsection{Top-p / nucleus sampling}

选择最小集合 $S_p$，使累计概率超过阈值 $p$：

\begin{equation}
S_p=\min\left\{S:
\sum_{i\in S}p_i\ge p
\right\}.
\end{equation}

归一化：

\begin{equation}
\tilde{p}_i=
\begin{cases}
\frac{p_i}{\sum_{j\in S_p}p_j}, & i\in S_p,\\
0, & i\notin S_p.
\end{cases}
\end{equation}

\subsection{Beam search}

维护 $K$ 条候选序列，序列分数：

\begin{equation}
\operatorname{score}(y_{1:T})
=\sum_{t=1}^{T}\log p_\theta(y_t\mid x,y_{<t}).
\end{equation}

常加长度惩罚：

\begin{equation}
\operatorname{score}_{\mathrm{lp}}(y)
=\frac{1}{T^\alpha}
\sum_{t=1}^{T}\log p_\theta(y_t\mid x,y_{<t}).
\end{equation}

\section{KV Cache 与解码复杂度}

\subsection{没有 KV cache 的代价}

生成第 $t$ 个 token 时，如果每一步都重算整个前缀，attention 总成本约为：

\begin{equation}
\sum_{t=1}^{T}O(t^2d)
=O(T^3d).
\end{equation}

实际 decoder 如果每步输入完整 prefix，会重复计算历史 K/V。

\subsection{有 KV cache 的增量解码}

第 $\ell$ 层缓存历史：

\begin{equation}
K_{\le t}^{\ell}
=\operatorname{Concat}(K_{\le t-1}^{\ell},K_t^{\ell}),
\quad
V_{\le t}^{\ell}
=\operatorname{Concat}(V_{\le t-1}^{\ell},V_t^{\ell}).
\end{equation}

新 token 只计算 query：

\begin{equation}
\bm{o}_t^\ell
=\softmax\left(
\frac{\bm{q}_t^\ell (K_{\le t}^{\ell})^\top}{\sqrt{d_h}}
\right)V_{\le t}^{\ell}.
\label{eq:kv-cache-attn}
\end{equation}

每步 attention 约 $O(td)$，整段生成约：

\begin{equation}
\sum_{t=1}^{T}O(td)=O(T^2d).
\end{equation}

\subsection{KV cache 显存公式}

假设 dtype bytes 为 $b$，层数 $L$，batch 中总上下文 token 数 $T_{\mathrm{tot}}$：

\begin{equation}
M_{\mathrm{KV}}
=2\cdot L\cdot T_{\mathrm{tot}}\cdot H_{\mathrm{kv}}\cdot d_h\cdot b.
\label{eq:kv-memory}
\end{equation}

其中 2 表示 K 和 V。对比：

\begin{equation}
M_{\mathrm{KV}}^{\mathrm{MHA}}:M_{\mathrm{KV}}^{\mathrm{GQA}}:M_{\mathrm{KV}}^{\mathrm{MQA}}
=H:G:1.
\end{equation}

\subsection{PagedAttention}

PagedAttention 把 KV cache 分成固定大小 blocks，而不是为每个 request 分配连续大块内存。设 block size 为 $B_{\mathrm{blk}}$，序列 $r$ 的 token 数为 $T_r$，需要 block 数：

\begin{equation}
n_r=\left\lceil\frac{T_r}{B_{\mathrm{blk}}}\right\rceil.
\end{equation}

逻辑位置 $t$ 映射到：

\begin{equation}
\operatorname{block}(t)=\left\lfloor\frac{t}{B_{\mathrm{blk}}}\right\rfloor,
\quad
\operatorname{offset}(t)=t\bmod B_{\mathrm{blk}}.
\end{equation}

这样可以减少 KV cache 碎片，并支持 prefix sharing/copy-on-write 等。来源：vLLM/PagedAttention \cite{kwon2023vllm}。

关系：

\[
\text{KV cache}
\rightarrow
\text{显存成为核心瓶颈}
\rightarrow
\text{PagedAttention/Block manager}
\rightarrow
\text{continuous batching 和高吞吐服务}.
\]

\section{Speculative Decoding}

Speculative decoding 使用小 draft model $q$ 先生成候选，再由大 target model $p$ 验证。

对 draft token $x$，接受概率：

\begin{equation}
\alpha(x)
=\min\left(1,\frac{p(x)}{q(x)}\right).
\label{eq:speculative-accept}
\end{equation}

若拒绝，则从修正分布采样：

\begin{equation}
p'(x)
\propto
\max(0,p(x)-q(x)).
\end{equation}

来源：Speculative Decoding \cite{leviathan2023speculative}。

直观关系：

\[
\text{draft model 快速提出 token}
\rightarrow
\text{target model 并行验证多个 token}
\rightarrow
\text{减少大模型解码步数}.
\]

\section{量化公式}

\subsection{Affine quantization}

把浮点 $x$ 量化为整数 $q$：

\begin{equation}
q=\clip\left(\operatorname{round}\left(\frac{x}{s}\right)+z,\ q_{\min},q_{\max}\right).
\end{equation}

反量化：

\begin{equation}
\hat{x}=s(q-z).
\end{equation}

其中 $s$ 是 scale，$z$ 是 zero point。

\subsection{Symmetric quantization}

对称量化常设 $z=0$：

\begin{equation}
s=\frac{\max |x|}{2^{b-1}-1},
\quad
q=\clip\left(\operatorname{round}\left(\frac{x}{s}\right),
-2^{b-1},2^{b-1}-1\right).
\end{equation}

来源可参考 LLM.int8 \cite{dettmers2022llmint8}、GPTQ \cite{frantar2022gptq}。

关系：

\[
\text{低 bit}
\rightarrow
\text{显存/带宽下降}
\rightarrow
\text{反量化或量化 kernel}
\rightarrow
\text{吞吐提升但精度可能下降}.
\]

\section{分布式训练/推理的核心公式}

\subsection{Data Parallel}

$R$ 个 worker 各算一份梯度：

\begin{equation}
\bm{g}
=\frac{1}{R}\sum_{r=1}^{R}\bm{g}^{(r)}.
\end{equation}

然后每个 worker 使用同一个平均梯度更新参数。

\subsection{Tensor Parallel：列切分}

线性层 $Y=XW$，按列切：

\begin{equation}
W=[W_1,W_2,\dots,W_R].
\end{equation}

每张卡计算：

\begin{equation}
Y_r=XW_r.
\end{equation}

最后 all-gather：

\begin{equation}
Y=[Y_1,Y_2,\dots,Y_R].
\end{equation}

\subsection{Tensor Parallel：行切分}

按行切：

\begin{equation}
X=[X_1,\dots,X_R],
\quad
W=
\begin{bmatrix}
W_1\\
\vdots\\
W_R
\end{bmatrix}.
\end{equation}

每张卡计算局部结果：

\begin{equation}
Y_r=X_rW_r.
\end{equation}

最后 all-reduce：

\begin{equation}
Y=\sum_{r=1}^{R}Y_r.
\end{equation}

\subsection{Expert Parallel}

MoE 中 expert 分布在不同 GPU。第 $t$ 个 token 被 router 分配到 expert $e_t$：

\begin{equation}
e_t=\argmax_e p_{t,e}.
\end{equation}

系统需要执行：

\[
\text{token dispatch}
\rightarrow
\text{expert local compute}
\rightarrow
\text{token combine}.
\]

数学上是式 \eqref{eq:moe}，系统上是 all-to-all communication。

\section{Multimodal 与 Diffusion 相关公式}

这一节不是纯 LLM 必需，但对理解 vLLM-Omni、Any-to-Any multimodal models 很重要。

\subsection{Vision patch embedding}

图像 $I\in\R^{H\times W\times C}$ 切成 $N$ 个 patch，每个 patch 展平为 $\bm{p}_i$：

\begin{equation}
\bm{v}_i=W_p\bm{p}_i+\bm{b}_p.
\end{equation}

然后把 image tokens 和 text tokens 拼接：

\begin{equation}
X=[X_{\mathrm{text}};X_{\mathrm{image}}].
\end{equation}

\subsection{CLIP 对比学习损失}

图文 embedding：

\begin{equation}
\bm{u}_i=f_{\mathrm{text}}(t_i),
\quad
\bm{v}_i=f_{\mathrm{image}}(I_i).
\end{equation}

相似度：

\begin{equation}
s_{ij}=\frac{\bm{u}_i^\top\bm{v}_j}{\tau}.
\end{equation}

image-to-text loss：

\begin{equation}
\mathcal{L}_{i2t}
=-\frac{1}{B}\sum_i
\log
\frac{\exp(s_{ii})}{\sum_j\exp(s_{ij})}.
\end{equation}

text-to-image 同理，最终：

\begin{equation}
\mathcal{L}_{\mathrm{CLIP}}
=\frac{1}{2}(\mathcal{L}_{i2t}+\mathcal{L}_{t2i}).
\end{equation}

来源：CLIP \cite{radford2021clip}。

\subsection{DDPM forward process}

扩散模型前向加噪：

\begin{equation}
q(\bm{x}_t\mid \bm{x}_{t-1})
=\mathcal{N}(\sqrt{1-\beta_t}\bm{x}_{t-1},\beta_t I).
\end{equation}

令 $\alpha_t=1-\beta_t$，$\bar{\alpha}_t=\prod_{s=1}^{t}\alpha_s$，可直接采样：

\begin{equation}
q(\bm{x}_t\mid \bm{x}_0)
=\mathcal{N}(\sqrt{\bar{\alpha}_t}\bm{x}_0,(1-\bar{\alpha}_t)I).
\end{equation}

即：

\begin{equation}
\bm{x}_t
=\sqrt{\bar{\alpha}_t}\bm{x}_0
+\sqrt{1-\bar{\alpha}_t}\bm{\epsilon},
\quad
\bm{\epsilon}\sim\mathcal{N}(0,I).
\end{equation}

来源：DDPM \cite{ho2020ddpm}。

\subsection{噪声预测训练目标}

模型预测噪声：

\begin{equation}
\mathcal{L}_{\mathrm{diff}}
=\E_{t,\bm{x}_0,\bm{\epsilon}}
\left[
\left\|
\bm{\epsilon}
-\bm{\epsilon}_\theta(\bm{x}_t,t,c)
\right\|_2^2
\right].
\label{eq:diffusion-loss}
\end{equation}

DiT 把 Transformer 用作 diffusion backbone。来源：DiT \cite{peebles2022dit}。

\subsection{Classifier-Free Guidance}

有条件预测 $\epsilon_\theta(\bm{x}_t,t,c)$ 和无条件预测 $\epsilon_\theta(\bm{x}_t,t,\varnothing)$：

\begin{equation}
\hat{\epsilon}
=\epsilon_\theta(\bm{x}_t,t,\varnothing)
+w\left[
\epsilon_\theta(\bm{x}_t,t,c)
-\epsilon_\theta(\bm{x}_t,t,\varnothing)
\right].
\label{eq:cfg}
\end{equation}

来源：Classifier-Free Guidance \cite{hosalimans2022cfg}。

和 vLLM-Omni 的关系：

\begin{itemize}
  \item AR stage 可产生 text/hidden/KV 条件。
  \item DiT stage 消费这些条件生成 image/video/audio。
  \item CFG companion requests 和 KV transfer 是系统实现层面的对应物。
\end{itemize}

\section{总关联图}

\subsection{训练侧}

\[
\begin{aligned}
\text{raw text}
&\xrightarrow{\text{tokenizer}}
x_{1:T}
\xrightarrow{E}
X\\
&\xrightarrow{\text{Transformer blocks}}
\bm{h}_{1:T}^L
\xrightarrow{W_U}
\bm{z}_{1:T}
\xrightarrow{\softmax}
p_\theta(x_{t+1}\mid x_{\le t})\\
&\xrightarrow{\mathrm{CE/NLL}}
\mathcal{L}
\xrightarrow{\mathrm{backprop}}
\nabla_\theta\mathcal{L}
\xrightarrow{\mathrm{AdamW}}
\theta_{t+1}.
\end{aligned}
\]

\subsection{推理侧}

\[
\begin{aligned}
\text{prompt}
&\xrightarrow{\text{prefill}}
K_{\le T},V_{\le T},\bm{h}_T\\
&\xrightarrow{\text{decode step}}
\bm{z}_{T}
\xrightarrow{\text{temperature/top-k/top-p}}
x_{T+1}\\
&\xrightarrow{\text{append KV cache}}
K_{\le T+1},V_{\le T+1}
\xrightarrow{\text{repeat}}
\text{generated text}.
\end{aligned}
\]

\subsection{系统侧}

\[
\begin{aligned}
\text{请求队列}
&\rightarrow
\text{batching/scheduling}
\rightarrow
\text{KV cache block allocation}\\
&\rightarrow
\text{GPU kernels: attention/FFN/sampling}
\rightarrow
\text{streaming output}.
\end{aligned}
\]

vLLM/SGLang 主要优化的是这一层，而不是改变语言建模目标本身。

\section{公式学习建议}

建议按下面顺序掌握：

\begin{enumerate}
  \item 先吃透式 \eqref{eq:ar-factorization}、\eqref{eq:nll}、\eqref{eq:softmax}、\eqref{eq:ce-gradient}。
  \item 再看 Transformer block：式 \eqref{eq:preln-block}、\eqref{eq:scaled-dot-product-attention}、\eqref{eq:causal-attention}。
  \item 再看现代模型变化：RoPE 式 \eqref{eq:rope-relative}、RMSNorm 式 \eqref{eq:rmsnorm}、SwiGLU 式 \eqref{eq:swiglu}。
  \item 然后看训练工程：AdamW 式 \eqref{eq:adamw}、Scaling Laws 式 \eqref{eq:scaling-laws}、计算量式 \eqref{eq:training-compute}。
  \item 微调对齐：SFT 式 \eqref{eq:sft}、Reward Model 式 \eqref{eq:reward-model}、PPO 式 \eqref{eq:ppo}、DPO 式 \eqref{eq:dpo}。
  \item 推理系统：KV cache 式 \eqref{eq:kv-cache-attn}、显存式 \eqref{eq:kv-memory}、PagedAttention block mapping。
  \item vLLM-Omni：再补 MoE 式 \eqref{eq:moe}、diffusion loss 式 \eqref{eq:diffusion-loss}、CFG 式 \eqref{eq:cfg}。
\end{enumerate}

\newpage
\begin{thebibliography}{99}

\bibitem{vaswani2017attention}
Ashish Vaswani et al.
\newblock Attention Is All You Need.
\newblock NeurIPS 2017.
\newblock \url{https://arxiv.org/abs/1706.03762}

\bibitem{devlin2018bert}
Jacob Devlin et al.
\newblock BERT: Pre-training of Deep Bidirectional Transformers for Language Understanding.
\newblock 2018.
\newblock \url{https://arxiv.org/abs/1810.04805}

\bibitem{radford2019gpt2}
Alec Radford et al.
\newblock Language Models are Unsupervised Multitask Learners.
\newblock OpenAI, 2019.
\newblock \url{https://cdn.openai.com/better-language-models/language_models_are_unsupervised_multitask_learners.pdf}

\bibitem{brown2020gpt3}
Tom Brown et al.
\newblock Language Models are Few-Shot Learners.
\newblock NeurIPS 2020.
\newblock \url{https://arxiv.org/abs/2005.14165}

\bibitem{sennrich2015bpe}
Rico Sennrich, Barry Haddow, Alexandra Birch.
\newblock Neural Machine Translation of Rare Words with Subword Units.
\newblock 2015.
\newblock \url{https://arxiv.org/abs/1508.07909}

\bibitem{kudo2018sentencepiece}
Taku Kudo, John Richardson.
\newblock SentencePiece: A simple and language independent subword tokenizer and detokenizer for Neural Text Processing.
\newblock 2018.
\newblock \url{https://aclanthology.org/D18-2012/}

\bibitem{kudo2018subword}
Taku Kudo.
\newblock Subword Regularization: Improving Neural Network Translation Models with Multiple Subword Candidates.
\newblock 2018.
\newblock \url{https://arxiv.org/abs/1804.10959}

\bibitem{ba2016layernorm}
Jimmy Lei Ba, Jamie Ryan Kiros, Geoffrey E. Hinton.
\newblock Layer Normalization.
\newblock 2016.
\newblock \url{https://arxiv.org/abs/1607.06450}

\bibitem{zhang2019rmsnorm}
Biao Zhang, Rico Sennrich.
\newblock Root Mean Square Layer Normalization.
\newblock 2019.
\newblock \url{https://arxiv.org/abs/1910.07467}

\bibitem{shazeer2020glu}
Noam Shazeer.
\newblock GLU Variants Improve Transformer.
\newblock 2020.
\newblock \url{https://arxiv.org/abs/2002.05202}

\bibitem{su2021rope}
Jianlin Su et al.
\newblock RoFormer: Enhanced Transformer with Rotary Position Embedding.
\newblock 2021.
\newblock \url{https://arxiv.org/abs/2104.09864}

\bibitem{press2021alibi}
Ofir Press, Noah A. Smith, Mike Lewis.
\newblock Train Short, Test Long: Attention with Linear Biases Enables Input Length Extrapolation.
\newblock 2021.
\newblock \url{https://arxiv.org/abs/2108.12409}

\bibitem{dao2022flashattention}
Tri Dao et al.
\newblock FlashAttention: Fast and Memory-Efficient Exact Attention with IO-Awareness.
\newblock 2022.
\newblock \url{https://arxiv.org/abs/2205.14135}

\bibitem{dao2023flashattention2}
Tri Dao.
\newblock FlashAttention-2: Faster Attention with Better Parallelism and Work Partitioning.
\newblock 2023.
\newblock \url{https://arxiv.org/abs/2307.08691}

\bibitem{shazeer2019mqa}
Noam Shazeer.
\newblock Fast Transformer Decoding: One Write-Head is All You Need.
\newblock 2019.
\newblock \url{https://arxiv.org/abs/1911.02150}

\bibitem{ainslie2023gqa}
Joshua Ainslie et al.
\newblock GQA: Training Generalized Multi-Query Transformer Models from Multi-Head Checkpoints.
\newblock 2023.
\newblock \url{https://arxiv.org/abs/2305.13245}

\bibitem{loshchilov2017adamw}
Ilya Loshchilov, Frank Hutter.
\newblock Decoupled Weight Decay Regularization.
\newblock 2017.
\newblock \url{https://arxiv.org/abs/1711.05101}

\bibitem{kaplan2020scaling}
Jared Kaplan et al.
\newblock Scaling Laws for Neural Language Models.
\newblock 2020.
\newblock \url{https://arxiv.org/abs/2001.08361}

\bibitem{hoffmann2022chinchilla}
Jordan Hoffmann et al.
\newblock Training Compute-Optimal Large Language Models.
\newblock 2022.
\newblock \url{https://arxiv.org/abs/2203.15556}

\bibitem{ouyang2022instructgpt}
Long Ouyang et al.
\newblock Training language models to follow instructions with human feedback.
\newblock 2022.
\newblock \url{https://arxiv.org/abs/2203.02155}

\bibitem{schulman2017ppo}
John Schulman et al.
\newblock Proximal Policy Optimization Algorithms.
\newblock 2017.
\newblock \url{https://arxiv.org/abs/1707.06347}

\bibitem{rafailov2023dpo}
Rafael Rafailov et al.
\newblock Direct Preference Optimization: Your Language Model is Secretly a Reward Model.
\newblock 2023.
\newblock \url{https://arxiv.org/abs/2305.18290}

\bibitem{hu2021lora}
Edward J. Hu et al.
\newblock LoRA: Low-Rank Adaptation of Large Language Models.
\newblock 2021.
\newblock \url{https://arxiv.org/abs/2106.09685}

\bibitem{lepikhin2020gshard}
Dmitry Lepikhin et al.
\newblock GShard: Scaling Giant Models with Conditional Computation and Automatic Sharding.
\newblock 2020.
\newblock \url{https://arxiv.org/abs/2006.16668}

\bibitem{fedus2021switch}
William Fedus, Barret Zoph, Noam Shazeer.
\newblock Switch Transformers: Scaling to Trillion Parameter Models with Simple and Efficient Sparsity.
\newblock 2021.
\newblock \url{https://arxiv.org/abs/2101.03961}

\bibitem{kwon2023vllm}
Woosuk Kwon et al.
\newblock Efficient Memory Management for Large Language Model Serving with PagedAttention.
\newblock 2023.
\newblock \url{https://arxiv.org/abs/2309.06180}

\bibitem{leviathan2023speculative}
Yaniv Leviathan, Matan Kalman, Yossi Matias.
\newblock Fast Inference from Transformers via Speculative Decoding.
\newblock 2023.
\newblock \url{https://arxiv.org/abs/2211.17192}

\bibitem{dettmers2022llmint8}
Tim Dettmers et al.
\newblock LLM.int8(): 8-bit Matrix Multiplication for Transformers at Scale.
\newblock 2022.
\newblock \url{https://arxiv.org/abs/2208.07339}

\bibitem{frantar2022gptq}
Elias Frantar et al.
\newblock GPTQ: Accurate Post-Training Quantization for Generative Pre-trained Transformers.
\newblock 2022.
\newblock \url{https://arxiv.org/abs/2210.17323}

\bibitem{radford2021clip}
Alec Radford et al.
\newblock Learning Transferable Visual Models From Natural Language Supervision.
\newblock 2021.
\newblock \url{https://arxiv.org/abs/2103.00020}

\bibitem{ho2020ddpm}
Jonathan Ho, Ajay Jain, Pieter Abbeel.
\newblock Denoising Diffusion Probabilistic Models.
\newblock 2020.
\newblock \url{https://arxiv.org/abs/2006.11239}

\bibitem{peebles2022dit}
William Peebles, Saining Xie.
\newblock Scalable Diffusion Models with Transformers.
\newblock 2022.
\newblock \url{https://arxiv.org/abs/2212.09748}

\bibitem{hosalimans2022cfg}
Jonathan Ho, Tim Salimans.
\newblock Classifier-Free Diffusion Guidance.
\newblock 2022.
\newblock \url{https://arxiv.org/abs/2207.12598}

\end{thebibliography}

\end{document}
